% ============================================================================
% CHAPTER 6: PRE-TRAINING, FINE-TUNING, AND TRANSFER LEARNING
% The practical backbone of modern NLP: how to adapt foundation models.
% ============================================================================

\chapter{Pre-training, Fine-tuning, and Transfer Learning}
\label{chap:pretraining_finetuning}

\begin{tldr}
The pre-train then fine-tune paradigm is the backbone of modern NLP.  This chapter
covers pre-training objectives (MLM, CLM, span corruption, replaced token detection,
denoising, contrastive), the pre-training data pipeline, full fine-tuning with
catastrophic forgetting mitigation, and the critical PEFT landscape---LoRA (the most
frequently asked topic), QLoRA, prefix tuning, adapters, and prompt tuning.
Understanding \emph{when} and \emph{how} to adapt foundation models is one of the
most practically relevant interview topics: every production NLP system built today
either fine-tunes a pre-trained model or uses one directly, and interviewers want to
know that you can reason about the full spectrum from pre-training data quality
through parameter-efficient adaptation to supervised fine-tuning for alignment.
\end{tldr}

\bigskip

% ============================================================================
% SECTION 1: PRE-TRAINING OBJECTIVES
% ============================================================================
\section{Pre-training Objectives}
\label{sec:pretraining-objectives}

Pre-training teaches a model general language representations from massive
unlabeled corpora before any task-specific supervision.  The choice of
pre-training objective fundamentally shapes what the model is good at---
understanding, generation, or both---and interviewers expect you to know
several objectives in detail, not just the one used by BERT.

\subsection{Masked Language Modeling (MLM)}

MLM is the pre-training objective introduced by BERT (Devlin et al., NAACL
2019).  The input sequence has tokens randomly selected for prediction, and
the model uses \emph{bidirectional} context to reconstruct them.

\paragraph{Masking procedure.}
15\% of input tokens are selected.  Of those:
\begin{itemize}
    \item \textbf{80\% replaced with \texttt{[MASK]}}---the primary training
          signal.
    \item \textbf{10\% replaced with a random token}---forces the model to
          maintain distributional representations for \emph{all} positions, not
          only masked ones.  Without this, the model would learn to ignore
          unmasked tokens entirely.
    \item \textbf{10\% kept unchanged}---biases the model toward the actual
          observed token distribution, reducing the train-inference mismatch
          (since \texttt{[MASK]} never appears during fine-tuning).
\end{itemize}

\begin{mathresult}{Masked Language Modeling Loss}
\[
    \loss_{\text{MLM}} = - \sum_{i \in \mathcal{M}} \log P\bigl(x_i \mid \vx_{\setminus \mathcal{M}}\bigr)
\]
where $\mathcal{M}$ is the set of masked positions and $\vx_{\setminus \mathcal{M}}$
denotes the corrupted input with masked, random, or unchanged tokens at positions in
$\mathcal{M}$.
\end{mathresult}

\paragraph{Why bidirectional context matters.}
Unlike autoregressive models that condition only on left context, MLM conditions
on \emph{both} left and right context simultaneously.  For understanding tasks
(classification, NER, extractive QA), this is strictly more informative:
``The cat sat on the \texttt{[MASK]}'' vs.\ ``The cat sat on the
\texttt{[MASK]} and purred'' yield very different predictions, and only
bidirectional models can leverage the right context.

\paragraph{Limitation.}
Only 15\% of tokens provide a training signal per forward pass---the other
85\% contribute to the contextual representation but generate no loss.  This
makes MLM sample-inefficient relative to objectives that provide signal at
every position.

\subsection{Causal Language Modeling (CLM)}

CLM is the objective used by all GPT-family models and the dominant LLM
pre-training approach.  The model predicts each token given only the
preceding tokens:

\begin{mathresult}{Causal Language Modeling Loss}
\[
    \loss_{\text{CLM}} = - \sum_{t=1}^{T} \log P\bigl(x_t \mid x_1, x_2, \ldots, x_{t-1}\bigr)
\]
\end{mathresult}

\paragraph{Why autoregressive enables generation.}
CLM directly models the joint distribution $P(\vx) = \prod_t P(x_t \mid x_{<t})$
via the chain rule of probability.  This factorization makes generation trivial:
sample one token at a time from the conditional distribution.  Bidirectional
models cannot do this because they condition on context that does not yet exist
during generation.

\paragraph{Every token provides signal.}
Unlike MLM where only 15\% of tokens are predicted, CLM produces a loss at
\emph{every} position, making it significantly more data-efficient per token
processed.  This is a key reason CLM scales better to massive corpora.

\paragraph{Limitation.}
Unidirectional context means the model cannot condition on future tokens for
understanding tasks, which is why encoder-only models (MLM) often outperform
similarly sized decoder-only models (CLM) on classification benchmarks when
both are fine-tuned.

\subsection{Span Corruption (T5-style)}

T5 (Raffel et al., JMLR 2020) introduced span corruption: randomly mask
contiguous spans of varying length, replace each with a single sentinel token,
and train the model to generate the masked spans in sequence.

\paragraph{Procedure.}
\begin{enumerate}
    \item Select approximately 15\% of tokens, grouped into spans (average
          span length $\approx 3$).
    \item Replace each span with a unique sentinel token (\texttt{<extra\_id\_0>},
          \texttt{<extra\_id\_1>}, etc.).
    \item The target sequence consists of sentinel tokens followed by the
          original text of each span.
\end{enumerate}

\paragraph{Advantages.}
Span masking is more natural than single-token masking because it forces the
model to generate coherent multi-token completions.  It works with an
encoder-decoder architecture where the encoder processes the corrupted input
bidirectionally and the decoder generates the spans autoregressively.  This
naturally supports both understanding (encoder) and generation (decoder).

\subsection{Replaced Token Detection (ELECTRA)}

ELECTRA (Clark et al., ICLR 2020) replaces MLM with a discriminative objective
that provides training signal at \emph{every} token position:

\paragraph{Architecture.}
\begin{enumerate}
    \item A small \textbf{generator} (trained with MLM) produces plausible
          replacement tokens for masked positions.
    \item The main \textbf{discriminator} receives the corrupted sequence
          (with generator outputs replacing masked tokens) and classifies
          each token as ``real'' or ``replaced.''
\end{enumerate}

\begin{mathresult}{ELECTRA Replaced Token Detection Loss}
\[
    \loss_{\text{RTD}} = - \sum_{t=1}^{T} \Bigl[
        y_t \log D(\tilde{x}_t) + (1 - y_t) \log\bigl(1 - D(\tilde{x}_t)\bigr)
    \Bigr]
\]
where $y_t = \mathbbm{1}[\tilde{x}_t = x_t]$ indicates whether position $t$ contains
the original token, $D$ is the discriminator, and $\tilde{\vx}$ is the sequence with
generator replacements.
\end{mathresult}

\begin{keyinsight}[Why ELECTRA Is More Efficient Than MLM]
MLM trains on 15\% of tokens per example.  ELECTRA trains on 100\%---every
token is classified as real or replaced.  This 6--7$\times$ increase in
training signal per example means ELECTRA matches BERT's performance with
roughly 1/4 the compute, and significantly outperforms BERT given the same
compute budget.  The key insight is that \emph{detection is easier than
generation} (binary classification vs.\ vocabulary-sized softmax), so even a
small discriminator can learn effectively from every position.
\end{keyinsight}

\subsection{Denoising Objectives (BART)}

BART (Lewis et al., ACL 2020) generalizes denoising to multiple corruption
strategies applied to an encoder-decoder architecture:

\begin{itemize}
    \item \textbf{Token masking:} Replace random tokens with \texttt{[MASK]}
          (similar to MLM).
    \item \textbf{Token deletion:} Remove random tokens without replacement
          (the model must figure out \emph{where} tokens are missing).
    \item \textbf{Text infilling:} Replace spans with a single \texttt{[MASK]}
          token (the model must determine span length).
    \item \textbf{Sentence permutation:} Shuffle sentence order within a document.
    \item \textbf{Document rotation:} Rotate the document to start at a random
          token.
\end{itemize}

The decoder must reconstruct the original, uncorrupted text.  Text infilling
was found to be the most effective single strategy, combining the benefits
of span masking with the challenge of length prediction.

\subsection{Contrastive Pre-training (for Embeddings)}

For models optimized to produce sentence or passage embeddings (used in
retrieval and RAG), contrastive objectives train the model to bring
semantically similar texts closer in embedding space while pushing
dissimilar texts apart.

\begin{mathresult}{InfoNCE Contrastive Loss}
\[
    \loss_{\text{InfoNCE}} = -\log \frac{\exp\bigl(\text{sim}(\vq, \vk^+) / \tau\bigr)}
        {\sum_{i=0}^{N} \exp\bigl(\text{sim}(\vq, \vk_i) / \tau\bigr)}
\]
where $\vq$ is the query representation, $\vk^+$ is the positive (matching)
key, $\{\vk_i\}$ includes the positive and $N$ negatives, $\tau$ is the
temperature, and $\text{sim}$ is typically cosine similarity.
\end{mathresult}

Models like E5, BGE, and GTE use contrastive pre-training (often on
large-scale web-mined text pairs) to produce high-quality embeddings
for retrieval, semantic search, and RAG.

\subsection{Comparison of Pre-training Objectives}

\begin{table}[H]
\centering
\small
\begin{tabularx}{\textwidth}{@{}l c c c X@{}}
\toprule
\textbf{Objective} & \textbf{Architecture} & \textbf{Signal/Token} & \textbf{Bidirectional?} & \textbf{Best For} \\
\midrule
MLM (BERT)          & Encoder-only   & 15\%   & Yes & Understanding, classification \\
CLM (GPT)           & Decoder-only   & 100\%  & No  & Generation, in-context learning \\
Span Corruption (T5)& Enc-Dec        & 15\%   & Encoder: Yes & Seq2seq tasks, translation \\
RTD (ELECTRA)       & Encoder-only   & 100\%  & Yes & Efficient understanding \\
Denoising (BART)    & Enc-Dec        & 100\%  & Encoder: Yes & Summarization, infilling \\
Contrastive         & Dual encoder   & 100\%  & Per-encoder & Retrieval, embeddings \\
\bottomrule
\end{tabularx}
\caption{Comparison of pre-training objectives.  ``Signal/Token'' indicates the fraction
of tokens that contribute a training signal per forward pass.}
\label{tab:pretraining-objectives}
\end{table}

\subsection{Visual Comparison: MLM vs.\ CLM vs.\ Span Corruption}

\begin{figure}[H]
\centering
\begin{tikzpicture}[
    node distance=0.3cm and 0.2cm,
    tok/.style={rectangle, draw=deepblue!50, fill=lightblue!60, minimum width=1.1cm,
                minimum height=0.5cm, font=\scriptsize\sffamily, rounded corners=1pt},
    masked/.style={rectangle, draw=deepred!60, fill=red!12, minimum width=1.1cm,
                   minimum height=0.5cm, font=\scriptsize\sffamily\bfseries, rounded corners=1pt},
    pred/.style={rectangle, draw=deepgreen!60, fill=green!10, minimum width=1.1cm,
                 minimum height=0.5cm, font=\scriptsize\sffamily\itshape, rounded corners=1pt},
    lbl/.style={font=\small\sffamily\bfseries, text=deepblue},
    arr/.style={-{Stealth[length=3pt]}, thick, gray!60}
]

% === MLM ===
\node[lbl] (mlm-lbl) at (0, 0) {MLM};
\node[tok, right=0.6cm of mlm-lbl] (m1) {The};
\node[tok, right=0.15cm of m1] (m2) {cat};
\node[masked, right=0.15cm of m2] (m3) {[MASK]};
\node[tok, right=0.15cm of m3] (m4) {on};
\node[tok, right=0.15cm of m4] (m5) {the};
\node[masked, right=0.15cm of m5] (m6) {[MASK]};
\node[pred, below=0.15cm of m3] (mp3) {sat};
\node[pred, below=0.15cm of m6] (mp6) {mat};
\draw[arr] (m3) -- (mp3);
\draw[arr] (m6) -- (mp6);

% === CLM ===
\node[lbl, below=1.2cm of mlm-lbl] (clm-lbl) {CLM};
\node[tok, right=0.6cm of clm-lbl] (c1) {The};
\node[tok, right=0.15cm of c1] (c2) {cat};
\node[tok, right=0.15cm of c2] (c3) {sat};
\node[tok, right=0.15cm of c3] (c4) {on};
\node[tok, right=0.15cm of c4] (c5) {the};
\node[pred, right=0.15cm of c5] (c6) {mat};
\node[pred, below=0.15cm of c2] (cp2) {cat};
\node[pred, below=0.15cm of c3] (cp3) {sat};
\node[pred, below=0.15cm of c4] (cp4) {on};
\node[pred, below=0.15cm of c5] (cp5) {the};
\draw[arr] (c1.south) -- ++(0,-0.15) -| (cp2.north);
\draw[arr] (c2.south) -- ++(0,-0.15) -| (cp3.north);
\draw[arr] (c3.south) -- ++(0,-0.15) -| (cp4.north);
\draw[arr] (c4.south) -- ++(0,-0.15) -| (cp5.north);

% === Span Corruption ===
\node[lbl, below=2.6cm of mlm-lbl] (span-lbl) {Span};
\node[tok, right=0.6cm of span-lbl] (s1) {The};
\node[masked, right=0.15cm of s1] (s2) {$\langle$X$\rangle$};
\node[tok, right=0.15cm of s2] (s3) {on};
\node[masked, right=0.15cm of s3] (s4) {$\langle$Y$\rangle$};
\node[pred, below=0.2cm of s2, xshift=0.6cm] (sp1) {$\langle$X$\rangle$\,cat sat};
\node[pred, below=0.2cm of s4, xshift=0.6cm] (sp2) {$\langle$Y$\rangle$\,the mat};
\draw[arr] (s2) -- (sp1);
\draw[arr] (s4) -- (sp2);

\end{tikzpicture}
\caption{Visual comparison of pre-training objectives.  \textbf{MLM}: predict masked
tokens from bidirectional context.  \textbf{CLM}: predict the next token from left
context at every position.  \textbf{Span Corruption}: replace variable-length spans
with sentinels, decode the original spans.}
\label{fig:pretraining-objectives}
\end{figure}


% ============================================================================
% SECTION 2: PRE-TRAINING DATA PIPELINE
% ============================================================================
\section{Pre-training Data Pipeline}
\label{sec:pretraining-data}

The quality and composition of pre-training data determine what a model knows,
what languages it speaks, and what biases it carries.  Interviewers increasingly
test data pipeline knowledge because it is the most impactful---and most
underappreciated---component of LLM development.

\subsection{Data Sources}

\begin{table}[H]
\centering
\small
\begin{tabularx}{\textwidth}{@{}l l X@{}}
\toprule
\textbf{Source} & \textbf{Scale} & \textbf{Characteristics} \\
\midrule
CommonCrawl / C4 / RefinedWeb & Trillions of tokens & Web text; noisy, diverse, broad coverage \\
Books (BookCorpus, Gutenberg) & Billions of tokens & Long-form, coherent prose; improves fluency \\
Wikipedia & $\sim$4B tokens (English) & Factual, well-structured, encyclopedic \\
Code (GitHub, StackOverflow) & Hundreds of billions & Improves reasoning and structured output \\
Scientific papers (ArXiv, S2ORC) & Tens of billions & Technical knowledge, formal writing \\
Curated instruction data & Millions of examples & For SFT stage; human-written or synthetic \\
\bottomrule
\end{tabularx}
\caption{Common pre-training data sources and their characteristics.}
\label{tab:data-sources}
\end{table}

\subsection{Data Quality Pipeline}

Raw web crawl data is far too noisy to use directly.  A typical cleaning pipeline:

\begin{figure}[H]
\centering
\begin{tikzpicture}[
    node distance=0.4cm,
    box/.style={rectangle, draw=deepblue, fill=lightblue!50, minimum width=3.0cm,
                minimum height=0.55cm, font=\scriptsize\sffamily, rounded corners=2pt, thick},
    arr/.style={-{Stealth[length=4pt]}, thick, deepblue!60}
]
\node[box] (raw) {Raw Web Crawl};
\node[box, below=0.4cm of raw] (lang) {Language Identification};
\node[box, below=0.4cm of lang] (dedup) {Deduplication (MinHash / Exact)};
\node[box, below=0.4cm of dedup] (quality) {Quality Filtering};
\node[box, below=0.4cm of quality] (toxic) {Toxicity / PII Filtering};
\node[box, below=0.4cm of toxic] (mix) {Domain Mixing};
\node[box, below=0.4cm of mix] (tokenize) {Tokenization};
\node[box, below=0.4cm of tokenize, fill=deepblue!15] (train) {Training Batches};

\draw[arr] (raw) -- (lang);
\draw[arr] (lang) -- (dedup);
\draw[arr] (dedup) -- (quality);
\draw[arr] (quality) -- (toxic);
\draw[arr] (toxic) -- (mix);
\draw[arr] (mix) -- (tokenize);
\draw[arr] (tokenize) -- (train);
\end{tikzpicture}
\caption{Typical pre-training data quality pipeline.  Each stage reduces data volume
while increasing quality.}
\label{fig:data-pipeline}
\end{figure}

\paragraph{Key stages.}
\begin{itemize}
    \item \textbf{Language identification:} Classify documents by language
          (fastText langid is standard).  Critical for multilingual models
          where per-language proportions are deliberately controlled.
    \item \textbf{Deduplication:} Remove near-duplicate documents using
          MinHash locality-sensitive hashing.  Exact substring deduplication
          catches repeated boilerplate (navigational text, legal notices).
          Lee et al.\ (2022) showed that deduplication improves both training
          efficiency and downstream quality.
    \item \textbf{Quality filtering:} Heuristic rules (minimum length,
          proportion of alphabetic characters, n-gram repetition ratio) and
          perplexity-based filtering (a language model scores documents;
          high-perplexity indicates gibberish or machine-generated noise).
    \item \textbf{Toxicity and PII filtering:} Classifiers remove toxic content;
          regex and NER-based systems redact personally identifiable information
          (names, emails, phone numbers, addresses).
\end{itemize}

\subsection{Data Mixing}

The proportion of each domain in the training set significantly affects model
capabilities:

\begin{itemize}
    \item \textbf{More code data} $\rightarrow$ better reasoning and structured
          output (an empirical finding replicated across Phi, CodeLLaMA, and
          DeepSeek models).
    \item \textbf{More book data} $\rightarrow$ better long-form coherence and
          narrative understanding.
    \item \textbf{More Wikipedia and scientific text} $\rightarrow$ better
          factual knowledge.
    \item \textbf{More multilingual data} $\rightarrow$ better cross-lingual
          transfer but potentially dilutes performance in dominant languages.
\end{itemize}

Determining optimal mixing proportions is an active research area.  The Chinchilla
paper (Hoffmann et al., 2022) established that compute-optimal training requires
approximately 20 tokens per parameter.  In practice, models are often
``over-trained'' (more tokens than Chinchilla-optimal) to reduce inference-time
model size while maximizing quality.

\subsection{Tokenizer Training}

The tokenizer is trained \emph{separately} from the model, before pre-training
begins.  Vocabulary size is a critical hyperparameter (typical range: 32K--100K
tokens).  The training corpus for the tokenizer determines what subwords emerge,
and this cannot be changed after model pre-training without retraining from
scratch.

\begin{warningbox}
Data quality is the single most underappreciated factor in LLM development.
The Phi model family demonstrated that small models trained on carefully curated
``textbook quality'' synthetic data can outperform much larger models trained on
raw web crawl.  In interviews, showing awareness that data quality often matters
more than model size is a strong signal of practical experience.
\end{warningbox}


% ============================================================================
% SECTION 3: TRANSFER LEARNING PARADIGM
% ============================================================================
\section{The Transfer Learning Paradigm}
\label{sec:transfer-learning}

Transfer learning---pre-train on a large, general corpus, then adapt to a
specific downstream task---is the fundamental paradigm of modern NLP.  It
replaced the previous approach of training task-specific models from scratch
with random initialization.

\subsection{Why Transfer Learning Works}

Pre-training learns \emph{universal language representations} that capture:
\begin{itemize}
    \item \textbf{Syntax:} Part-of-speech patterns, constituency structure,
          dependency relations.
    \item \textbf{Semantics:} Word meaning, synonymy, analogy, compositional
          meaning.
    \item \textbf{World knowledge:} Factual associations, commonsense reasoning,
          entity relationships.
    \item \textbf{Discourse:} Coherence, coreference, topic structure.
\end{itemize}

These representations transfer across tasks because language understanding has
common prerequisites: a model that understands syntax and semantics for MLM
also understands them for NER, QA, and sentiment analysis.  This dramatically
reduces the labeled data needed for any individual task.

\subsection{Feature Extraction vs.\ Fine-tuning}

\begin{table}[H]
\centering
\small
\begin{tabularx}{\textwidth}{@{}l X X@{}}
\toprule
\textbf{Aspect} & \textbf{Feature Extraction} & \textbf{Fine-tuning} \\
\midrule
What is trained & Only the task head & All parameters (or most) \\
Pre-trained weights & Frozen & Updated \\
Training cost & Very low & Moderate to high \\
Data requirement & Can work with very little data & Needs hundreds+ examples \\
Quality & Good baseline & Usually significantly better \\
Risk & No catastrophic forgetting & Possible catastrophic forgetting \\
Use case & Quick baseline, limited compute & Production systems \\
\bottomrule
\end{tabularx}
\caption{Feature extraction vs.\ fine-tuning for transfer learning.}
\label{tab:feature-vs-finetune}
\end{table}

\paragraph{Feature extraction.}
Freeze the pre-trained model, extract hidden-state representations for each
input, and train a lightweight classifier (linear layer, small MLP) on top.
The pre-trained model is used purely as a feature extractor.

\paragraph{Fine-tuning.}
Unfreeze all (or most) parameters and train the entire model end-to-end on
the downstream task.  This allows the pre-trained representations to adapt to
the task-specific distribution, typically yielding 2--5\% absolute improvement
over feature extraction.

\subsection{Transfer Learning Pipeline}

\begin{figure}[H]
\centering
\begin{tikzpicture}[
    node distance=0.8cm,
    phase/.style={rectangle, draw=deepblue, fill=lightblue!40, minimum width=3.5cm,
                  minimum height=0.8cm, font=\small\sffamily, rounded corners=3pt, thick},
    data/.style={rectangle, draw=gray!50, fill=gray!8, minimum width=3.5cm,
                 minimum height=0.5cm, font=\scriptsize\sffamily, rounded corners=2pt},
    arr/.style={-{Stealth[length=5pt]}, thick, deepblue!70},
    lbl/.style={font=\scriptsize\sffamily, text=gray!70}
]

\node[phase] (pretrain) {Pre-training};
\node[data, above=0.2cm of pretrain] (pd) {Unlabeled corpus (TB-scale)};
\node[lbl, below=0.05cm of pretrain] {MLM / CLM / Span Corruption};

\node[phase, right=2.5cm of pretrain] (finetune) {Fine-tuning};
\node[data, above=0.2cm of finetune] (fd) {Labeled task data (100s--10Ks)};
\node[lbl, below=0.05cm of finetune] {Task-specific loss};

\node[phase, right=2.5cm of finetune, fill=deepblue!15] (deploy) {Deployment};
\node[data, above=0.2cm of deploy] (dd) {Production queries};
\node[lbl, below=0.05cm of deploy] {Inference};

\draw[arr] (pretrain) -- (finetune) node[midway, above, font=\scriptsize\sffamily] {Transfer weights};
\draw[arr] (finetune) -- (deploy) node[midway, above, font=\scriptsize\sffamily] {Optimized model};

\end{tikzpicture}
\caption{The transfer learning pipeline: pre-train on a large unlabeled corpus,
fine-tune on task-specific labeled data, deploy for inference.}
\label{fig:transfer-pipeline}
\end{figure}

\subsection{Layer-wise Learning Rates}

Different layers capture different levels of abstraction.  Lower layers learn
general features (syntax, morphology) that transfer broadly; upper layers learn
more task-specific features.  A common fine-tuning strategy assigns
\emph{lower learning rates to earlier layers} and higher rates to later layers:

\[
    \eta_l = \eta_{\text{base}} \cdot \xi^{L - l}
\]

where $\eta_l$ is the learning rate for layer $l$, $L$ is the total number of
layers, and $\xi \in (0, 1)$ is the decay factor (typically 0.9--0.95).
Howard and Ruder (ULMFiT, ACL 2018) showed this discriminative fine-tuning
significantly outperforms uniform learning rates.


% ============================================================================
% SECTION 4: FULL FINE-TUNING
% ============================================================================
\section{Full Fine-Tuning}
\label{sec:full-finetuning}

Full fine-tuning updates \emph{all} model parameters on the downstream task.
It is the simplest and often most effective adaptation method when sufficient
data and compute are available.

\subsection{Standard Recipe}

\begin{enumerate}
    \item \textbf{Initialize} from pre-trained weights.
    \item \textbf{Add a task head:} linear layer(s) on top of the pre-trained
          model's output (e.g., a classification head on \texttt{[CLS]}, or
          token-level heads for NER).
    \item \textbf{Train all parameters} end-to-end with a task-specific loss
          (cross-entropy for classification, span extraction loss for QA).
    \item \textbf{Use a low learning rate:} typically $1 \times 10^{-5}$ to
          $5 \times 10^{-5}$ for the pre-trained backbone (the task head can
          use a higher rate).
    \item \textbf{Train for few epochs:} 2--5 epochs is usually sufficient for
          fine-tuning.  Unlike training from scratch, the model starts from a
          strong initialization.
\end{enumerate}

\subsection{Catastrophic Forgetting}
\label{sec:catastrophic-forgetting}

\textbf{Catastrophic forgetting} occurs when fine-tuning on a new task causes
the model to lose knowledge acquired during pre-training.  The parameters
shift to optimize the fine-tuning loss, potentially overwriting representations
that were useful for general language understanding.

\paragraph{Why it happens.}
Pre-training and fine-tuning optimize \emph{different} objectives on
\emph{different} data distributions.  Without explicit constraints, gradient
updates for the fine-tuning task can destructively interfere with the parameter
configurations that encoded pre-trained knowledge.  The smaller the fine-tuning
dataset, the more severe the problem---the model overfits to the narrow
distribution and loses generality.

\subsection{Mitigating Catastrophic Forgetting}

\begin{enumerate}
    \item \textbf{Low learning rate} ($1\text{e-}5$ to $5\text{e-}5$):
          Small updates preserve the pre-trained parameter neighborhood.
          This is the single most important mitigation.

    \item \textbf{Gradual unfreezing} (ULMFiT, Howard \& Ruder, 2018):
          Start by training only the task head, then progressively unfreeze
          layers from top to bottom.  This allows lower layers to remain
          stable while upper layers adapt first.

    \item \textbf{Regularization:}
          \begin{itemize}
              \item Weight decay ($L_2$ regularization) penalizes large parameter
                    changes from the pre-trained values.
              \item Dropout (typically 0.1) prevents co-adaptation.
              \item Elastic Weight Consolidation (EWC): adds a penalty
                    $\sum_i F_i (\theta_i - \theta_i^*)^2$ where $F_i$ is the
                    Fisher information of parameter $i$ and $\theta_i^*$ are the
                    pre-trained values.  Parameters important for pre-training
                    are penalized more for changing.
          \end{itemize}

    \item \textbf{Early stopping:} Monitor validation performance and stop when
          it degrades.  Fine-tuning often peaks quickly (2--3 epochs) before
          overfitting.

    \item \textbf{Data replay:} Mix a small proportion of pre-training-style
          data with fine-tuning data to maintain general representations.
\end{enumerate}

\begin{warningbox}
The most common fine-tuning mistakes in practice (and the ones interviewers
look for): (1) Using too high a learning rate (e.g., $1\text{e-}3$ instead of
$1\text{e-}5$), which destroys pre-trained representations in the first few
steps. (2) Training for too many epochs on a small dataset, causing severe
overfitting. (3) Not using a learning rate warmup, which can cause unstable
early updates. (4) Forgetting to set the model to training mode (missing
dropout, incorrect batch norm behavior). (5) Not freezing the embedding layer
when the task vocabulary matches the pre-trained tokenizer exactly.
\end{warningbox}


% ============================================================================
% SECTION 5: PARAMETER-EFFICIENT FINE-TUNING (PEFT)
% ============================================================================
\section{Parameter-Efficient Fine-Tuning (PEFT)}
\label{sec:peft}

As models grow to billions or hundreds of billions of parameters, full
fine-tuning becomes prohibitively expensive in memory, compute, and storage.
\textbf{Parameter-Efficient Fine-Tuning (PEFT)} methods train only a small
fraction of parameters while keeping the vast majority frozen, achieving
near-full-fine-tuning quality at a fraction of the cost.

\subsection{LoRA: Low-Rank Adaptation --- Deep Dive}
\label{sec:lora}

LoRA (Hu et al., ICLR 2022) is the most widely used PEFT method and the
single most frequently asked fine-tuning topic in NLP interviews.

\subsubsection{Core Hypothesis: Low-Rank Updates}

The key insight behind LoRA is that \emph{the weight updates during fine-tuning
occupy a low-rank subspace}.  Pre-trained models already have high-quality
representations; fine-tuning only needs to make small, structured adjustments
to adapt them.  Aghajanyan et al.\ (2021) formalized this as the
\textbf{intrinsic dimensionality} hypothesis: the optimization landscape of
fine-tuning can be effectively navigated in a subspace of dimension much
smaller than the full parameter space.

\subsubsection{Mathematical Formulation}

For a pre-trained weight matrix $\mW_0 \in \R^{d \times k}$, LoRA
parameterizes the update as a low-rank factorization:

\begin{mathresult}{LoRA: Low-Rank Adaptation}
\[
    \mW = \mW_0 + \Delta\mW = \mW_0 + \mB\mA
\]
where $\mB \in \R^{d \times r}$, $\mA \in \R^{r \times k}$, and $r \ll \min(d, k)$.

During training, $\mW_0$ is \textbf{frozen} (no gradient updates).
Only $\mA$ and $\mB$ are trained.  The forward pass becomes:
\[
    \vh = \mW_0 \vx + \mB\mA\vx = \mW_0 \vx + \Delta\mW\,\vx
\]

With scaling:
\[
    \vh = \mW_0 \vx + \frac{\alpha}{r}\,\mB\mA\vx
\]
where $\alpha$ is a constant scaling factor.
\end{mathresult}

\subsubsection{Initialization}

Correct initialization is critical for stable training:
\begin{itemize}
    \item $\mA$ is initialized with a random Gaussian: $A_{ij} \sim \mathcal{N}(0, \sigma^2)$.
    \item $\mB$ is initialized to \textbf{zero}: $\mB = \mathbf{0}$.
\end{itemize}
This ensures that $\Delta\mW = \mB\mA = \mathbf{0}$ at initialization, so the
model starts from exactly the pre-trained weights.  During training, $\mB$
gradually moves away from zero as gradients flow.  This zero initialization is
analogous to the residual connection design principle: start as identity, learn
the residual.

\subsubsection{Rank $r$: The Key Hyperparameter}

The rank $r$ controls the expressiveness of the update:

\begin{itemize}
    \item \textbf{Typical range:} $r = 8$ to $64$.
    \item \textbf{Low rank ($r = 4$--$8$):} Sufficient for tasks similar to the
          pre-training domain.  Extremely parameter-efficient.
    \item \textbf{Medium rank ($r = 16$--$32$):} Good default for most
          fine-tuning scenarios.
    \item \textbf{High rank ($r = 64$--$128$):} Approaches full fine-tuning
          capacity.  Diminishing returns beyond $r = 64$ for most tasks.
    \item \textbf{Rule of thumb:} More complex tasks or larger domain shifts
          require higher rank.  But the relationship is sublinear---doubling
          $r$ does not double quality.
\end{itemize}

\subsubsection{Alpha ($\alpha$) Scaling}

The effective magnitude of the LoRA update is:
\[
    \text{effective update} = \frac{\alpha}{r} \cdot \mB\mA
\]

The ratio $\alpha/r$ controls how strongly the adaptation modifies the original
weights.  Common practice:
\begin{itemize}
    \item Set $\alpha = 2r$ (so the effective scaling is 2) as a starting point.
    \item Higher $\alpha$ means stronger adaptation; lower means more conservative.
    \item When sweeping $r$, keeping $\alpha/r$ constant maintains similar
          update magnitude across different rank settings.
\end{itemize}

\subsubsection{Which Modules to Apply LoRA}

\begin{itemize}
    \item \textbf{Standard:} Apply to the query ($\mW_Q$) and value ($\mW_V$)
          projections in attention.  This is the default in the original paper.
    \item \textbf{Extended:} Also apply to the key ($\mW_K$) and output
          ($\mW_O$) projections for slightly better quality.
    \item \textbf{Full:} Apply to all linear layers including FFN projections
          ($\mW_1$, $\mW_2$, and gate projection for SwiGLU).  Higher quality
          but more parameters.
    \item \textbf{Empirical finding:} For most tasks, Q+V is sufficient.  Adding
          FFN projections helps when the task requires significant knowledge
          adaptation (e.g., domain-specific fine-tuning).
\end{itemize}

\subsubsection{Parameter Count}

For a single weight matrix $\mW_0 \in \R^{d \times k}$:
\begin{align*}
    \text{Full fine-tuning params} &= d \times k \\
    \text{LoRA params} &= d \times r + r \times k = r(d + k)
\end{align*}

For a 7B-parameter model with $d = 4096$, applying rank-16 LoRA to Q and V
projections across 32 layers:
\[
    \text{LoRA params} = 32 \times 2 \times 16 \times (4096 + 4096)
    = 8{,}388{,}608 \approx 8.4\text{M}
\]
That is approximately \textbf{0.12\%} of the full model parameters.

\subsubsection{Inference: Zero Overhead via Merging}

\begin{keyinsight}[LoRA Merging: Zero Inference Overhead]
After training, the LoRA matrices can be \emph{merged} into the base weights:
$\mW_{\text{merged}} = \mW_0 + (\alpha/r) \cdot \mB\mA$.  The merged model has
\emph{exactly the same architecture and inference cost} as the original---no
additional latency, no additional memory for LoRA matrices.  This is a
decisive advantage over methods like adapters or prefix tuning that add
inference overhead.  Furthermore, you can maintain multiple LoRA adapters
(one per task) and swap them at serving time by loading different $\mA, \mB$
matrices without reloading the base model.
\end{keyinsight}

\subsubsection{Multi-task LoRA Serving}

In production, a common pattern is to serve multiple tasks from a single base
model with different LoRA adapters:

\begin{enumerate}
    \item Load the base model once into GPU memory.
    \item For each task, load the corresponding $\mA, \mB$ matrices
          (typically a few MB per adapter).
    \item At request time, select the appropriate adapter based on the task
          and either merge on-the-fly or maintain pre-merged copies.
    \item This enables a single infrastructure to serve dozens of specialized
          tasks with minimal memory overhead.
\end{enumerate}

\subsubsection{Why Low-Rank Works: The Intrinsic Dimensionality Argument}

\begin{keyinsight}[Intrinsic Dimensionality of Fine-Tuning]
Aghajanyan et al.\ (2021) showed that the effective dimensionality of the
fine-tuning optimization landscape is orders of magnitude smaller than the
full parameter count.  For BERT-base (110M parameters), the intrinsic
dimension for many NLP tasks is as low as $\sim$200--800.  This means the
``useful'' directions for adaptation lie in a tiny subspace of the full
parameter space.  LoRA directly exploits this structure: rank $r = 16$ provides
a $16$-dimensional adaptation per weight matrix, and across all adapted matrices,
the total intrinsic dimensionality easily covers the effective subspace needed
for most tasks.
\end{keyinsight}

\subsection{QLoRA: Quantized LoRA}
\label{sec:qlora}

QLoRA (Dettmers et al., NeurIPS 2023) makes LoRA applicable to models that
would otherwise not fit in memory by combining quantization with low-rank
adaptation.

\subsubsection{Three Key Innovations}

\begin{enumerate}
    \item \textbf{4-bit NormalFloat (NF4) quantization:} The base model
          weights are stored in 4-bit precision using a data type optimized
          for the normal distribution of pre-trained weights.  NF4 quantization
          bins are spaced to minimize quantization error for normally
          distributed values, unlike uniform INT4.

    \item \textbf{Double quantization:} The quantization constants (scales
          and zero-points, one per block of 64 weights) are themselves quantized
          to 8-bit, reducing the memory overhead of quantization metadata from
          $\sim$0.5 bits/parameter to $\sim$0.13 bits/parameter.

    \item \textbf{Paged optimizers:} Use CPU RAM to handle GPU memory spikes
          during gradient computation via unified memory, preventing
          out-of-memory errors during training.
\end{enumerate}

\subsubsection{Computation Flow}

During the forward pass, 4-bit weights are \emph{dequantized} on-the-fly to
16-bit for matrix multiplication.  The LoRA adapters ($\mA, \mB$) are always
stored and computed in 16-bit (BF16 or FP16).  Gradients flow only through the
LoRA parameters---the quantized base weights are never updated.

\begin{mathresult}{QLoRA Memory Savings}
For a model with $N$ parameters:
\begin{align*}
    \text{Full fine-tuning memory} &\approx 16N \text{ bytes}
    \quad \text{(FP16 weights + FP32 optimizer states)} \\[4pt]
    \text{QLoRA memory} &\approx 0.5N + 16 \cdot r(d+k) \cdot L_{\text{adapted}}
    \text{ bytes}
\end{align*}
where the first term is the 4-bit base model ($\sim$0.5 bytes/parameter) and
the second term is the 16-bit LoRA parameters.  For a 65B model with rank-16 LoRA:
full fine-tuning requires $\sim$780\,GB; QLoRA requires $\sim$33\,GB (single
48GB GPU).
\end{mathresult}

\paragraph{Quality.}
QLoRA achieves 85--95\% of full fine-tuning quality.  The gap comes primarily
from the 4-bit quantization of the base model, not from the low-rank constraint
of LoRA.  For many practical tasks, this quality gap is negligible.

\subsection{DoRA: Decomposed Weight-Magnitude and Direction}
\label{sec:dora}

DoRA (Liu et al., 2024) improves on LoRA by decomposing the weight update into
\emph{magnitude} and \emph{direction} components:

\[
    \mW' = m \cdot \frac{\mW_0 + \mB\mA}{\norm{\mW_0 + \mB\mA}_c}
\]

where $m$ is a learnable magnitude vector and $\norm{\cdot}_c$ denotes
column-wise normalization.  The direction is updated via LoRA, and the magnitude
is updated independently.  This decomposition is inspired by weight normalization
(Salimans \& Kingma, 2016) and allows the model to adjust the ``how much''
and ``in which direction'' of the update separately.  DoRA consistently
outperforms LoRA by 0.5--1\% across benchmarks with the same parameter budget.

\subsection{Prefix Tuning}
\label{sec:prefix-tuning}

Prefix tuning (Li \& Liang, ACL 2021) prepends learnable continuous vectors
(``virtual tokens'') to the keys and values at every transformer layer:

\[
    \mK' = [\mK_{\text{prefix}};\; \mK], \qquad
    \mV' = [\mV_{\text{prefix}};\; \mV]
\]

where $\mK_{\text{prefix}} \in \R^{p \times d_k}$ and
$\mV_{\text{prefix}} \in \R^{p \times d_v}$ are $p$ learnable prefix vectors
per layer.  The actual input tokens attend to both the prefix and each other,
while the prefix is trained end-to-end.

\paragraph{Key properties.}
Trainable parameters: $2 \times p \times d \times L$ (prefix length $\times$
dimension $\times$ layers, for K and V).  Typically $p = 10$--$50$, yielding
$\sim$0.1\% of total parameters.  Quality: 85--95\% of full fine-tuning, best
when the task is well-defined and does not require substantial knowledge
adaptation.

\paragraph{Limitation.}
The prefix consumes context window positions, effectively reducing the usable
context length.  Also, prefix tuning adds inference overhead because the
prefix vectors must be processed at every layer for every forward pass.

\subsection{Prompt Tuning}
\label{sec:prompt-tuning}

Prompt tuning (Lester et al., EMNLP 2021) is the simplest PEFT method:
prepend $p$ learnable soft tokens to the input embedding only (not at every
layer, unlike prefix tuning):

\[
    \mX' = [\mE_{\text{prompt}};\; \mX]
\]

where $\mE_{\text{prompt}} \in \R^{p \times d}$ are learnable embeddings.

\paragraph{Key properties.}
Fewest trainable parameters of any PEFT method ($< 0.1\%$).  Performance
improves with model scale: for models under 1B parameters, prompt tuning
underperforms significantly; for models above 10B, it approaches full
fine-tuning quality.  The intuition: larger models can ``interpret'' soft
prompts more effectively because their in-context learning ability is
stronger.

\subsection{Adapters}
\label{sec:adapters}

Adapter layers (Houlsby et al., ICML 2019) insert small bottleneck modules
between existing transformer sublayers:

\[
    \vh = \vh + f(\vh\,\mW_{\text{down}})\,\mW_{\text{up}}
\]

where $\mW_{\text{down}} \in \R^{d \times m}$ projects to a bottleneck
dimension $m \ll d$, a nonlinearity $f$ (ReLU or GELU) is applied, and
$\mW_{\text{up}} \in \R^{m \times d}$ projects back.  A residual connection
ensures the adapter starts as identity.

\paragraph{Trainable parameters:} $\sim$1--5\% of the model, depending on
bottleneck size $m$.  Quality: 95--99\% of full fine-tuning.  Limitation:
unlike LoRA, adapters cannot be merged into the base weights, so they add
inference latency (two additional matrix multiplications per layer).

\subsection{IA3: Infused Adapter by Inhibiting and Amplifying Inner Activations}
\label{sec:ia3}

IA3 (Liu et al., 2022) learns three vectors $\mathbf{l}_k, \mathbf{l}_v,
\mathbf{l}_{ff}$ that \emph{rescale} the keys, values, and FFN activations
element-wise:

\[
    \vk' = \mathbf{l}_k \odot \vk, \qquad
    \mathbf{v}' = \mathbf{l}_v \odot \mathbf{v}, \qquad
    \vh'_{ff} = \mathbf{l}_{ff} \odot \vh_{ff}
\]

This is the most extreme form of parameter efficiency: only three vectors per
layer, yielding $\sim$0.01\% trainable parameters.  Quality: 85--90\% of full
fine-tuning, best for few-shot adaptation.

\subsection{Comprehensive PEFT Comparison}

\begin{table}[H]
\centering
\small
\begin{tabularx}{\textwidth}{@{}l l c c X@{}}
\toprule
\textbf{Method} & \textbf{Mechanism} & \textbf{\% Params} & \textbf{Quality vs.\ Full FT} & \textbf{Inference Overhead} \\
\midrule
LoRA          & Low-rank matrices on attn   & 0.1--1\%    & 90--99\%  & None (merge) \\
QLoRA         & LoRA on 4-bit base model    & 0.1--1\%    & 85--95\%  & None (merge) \\
DoRA          & Magnitude + direction LoRA  & 0.1--1\%    & 92--99\%  & None (merge) \\
Prefix Tuning & Virtual KV tokens per layer & $\sim$0.1\% & 85--95\%  & Extra KV at every layer \\
Prompt Tuning & Soft input tokens only      & $< 0.1$\%   & 80--90\%  & Extra input tokens \\
Adapters      & Bottleneck modules          & 1--5\%      & 95--99\%  & Two MatMuls per layer \\
IA3           & Learned scaling vectors     & $\sim$0.01\%& 85--90\%  & Negligible (element-wise) \\
\bottomrule
\end{tabularx}
\caption{Comprehensive comparison of PEFT methods.  ``Quality vs.\ Full FT'' indicates
the typical fraction of full fine-tuning performance retained.}
\label{tab:peft-comparison}
\end{table}


% ============================================================================
% SECTION 6: SUPERVISED FINE-TUNING (SFT) FOR LLMs
% ============================================================================
\section{Supervised Fine-Tuning (SFT) for LLMs}
\label{sec:sft}

Supervised fine-tuning transforms a pre-trained language model (which generates
arbitrary text completions) into an \emph{instruction-following assistant}
(which provides helpful, structured responses).  SFT is the bridge between
pre-training and alignment (RLHF/DPO, covered in Chapter~\ref{chap:pretraining_finetuning}+1).

\subsection{Data Format}

SFT data consists of instruction-response pairs, typically organized in a chat
format:

\begin{lstlisting}[language={}, caption={Typical SFT training example (ChatML format)}]
<|system|>
You are a helpful assistant.
<|user|>
Explain the difference between LoRA and full fine-tuning.
<|assistant|>
LoRA (Low-Rank Adaptation) freezes the pre-trained weights
and adds small trainable low-rank matrices...
\end{lstlisting}

Common chat templates include ChatML (used by OpenAI-style models), Alpaca
format (instruction/input/output fields), LLaMA chat format (with
\texttt{[INST]} tags), and Vicuna format.  The exact template
\emph{must match} what the model was trained with---using the wrong template
at inference produces degraded outputs.

\subsection{Loss Masking}

A critical implementation detail: during SFT, the loss is computed
\textbf{only on the assistant tokens}, not on the system prompt or user
message:

\begin{mathresult}{SFT Loss with Masking}
\[
    \loss_{\text{SFT}} = -\sum_{t \in \mathcal{A}} \log P(x_t \mid x_{<t})
\]
where $\mathcal{A}$ is the set of positions corresponding to the assistant
response.  User and system tokens provide context but generate no gradient signal.
\end{mathresult}

\paragraph{Why mask the prompt.}
The model should learn to \emph{generate good responses}, not to predict the
user's question (which would waste capacity learning to model the distribution
of user inputs rather than the conditional distribution of good responses given
inputs).

\subsection{Data Quality Matters More Than Quantity}

\begin{keyinsight}[The LIMA Principle: Less Is More for Alignment]
Zhou et al.\ (2023) demonstrated with LIMA that fine-tuning LLaMA-65B on just
\textbf{1,000 carefully curated} instruction-response pairs produced a model
competitive with GPT-4-era instruction-tuned models on many benchmarks.  The
key finding: a pre-trained model already has the knowledge; SFT mainly teaches
it the \emph{format and style} of helpful responses.  A few high-quality
examples are more effective than thousands of noisy ones because they provide a
clear, consistent signal about how to behave.
\end{keyinsight}

\paragraph{What makes SFT data high-quality.}
\begin{itemize}
    \item \textbf{Correctness:} Responses must be factually accurate.
    \item \textbf{Diversity:} Cover many task types, complexity levels, and
          formats (Q\&A, summarization, coding, math, creative writing, etc.).
    \item \textbf{Consistency:} All examples should follow the same style,
          formatting, and safety conventions.
    \item \textbf{Complexity:} Include genuinely hard examples, not just
          trivial ones---the model needs to learn how to handle difficulty.
    \item \textbf{Multi-turn coverage:} Include multi-turn conversations to
          teach context management, not just single-turn Q\&A.
\end{itemize}

\subsection{SFT as the Bridge to Alignment}

The three-stage LLM training pipeline is:
\[
    \text{Pre-training} \xrightarrow{\text{large unlabeled}} \text{SFT}
    \xrightarrow{\text{small labeled}} \text{RLHF/DPO}
    \xrightarrow{\text{preferences}} \text{Aligned Model}
\]

SFT provides the behavioral foundation: it teaches the model the
\emph{format} of being helpful (respond in complete sentences, follow
instructions, refuse harmful requests).  Alignment (RLHF/DPO) then refines
the \emph{quality} of that behavior using preference data.  Without SFT, RLHF
struggles because the model does not know how to produce well-formed responses
that can be meaningfully compared.


% ============================================================================
% SECTION 7: INTERVIEW QUESTIONS
% ============================================================================
\section{Interview Questions}
\label{sec:pretraining-questions}

%% QUESTION 1 %%
\begin{interviewq}
\textbf{Q: Explain LoRA. Why does it work? How do you choose rank and alpha?}

\badgeLfive{L5} \quad \badgetype{Mathematical} \quad \badgefreq{\freqhigh}

\stronganswer
LoRA freezes the pre-trained weight matrix $\mW_0 \in \R^{d \times k}$ and adds a
low-rank update $\Delta\mW = \mB\mA$ where $\mB \in \R^{d \times r}$ and
$\mA \in \R^{r \times k}$, with $r \ll \min(d, k)$.  Only $\mA$ and $\mB$ are
trained---typically 0.1--1\% of total parameters.  It works because fine-tuning
updates are empirically low-rank: the intrinsic dimensionality of the fine-tuning
optimization landscape is much smaller than the full parameter count (Aghajanyan
et al., 2021 showed intrinsic dimension $\sim$200--800 for BERT tasks on 110M
parameters).  The model already has good representations; adaptation only needs
adjustments in a low-dimensional subspace.

$\mB$ is initialized to zero and $\mA$ to random Gaussian so that $\Delta\mW = 0$
at the start---the model begins from exactly the pre-trained weights.  The effective
update is scaled by $\alpha/r$, where $\alpha$ controls adaptation strength.  For
rank: $r = 16$--$32$ is a good default; increase for complex tasks or large domain
shifts.  For alpha: $\alpha = 2r$ is a common starting point.  Keeping $\alpha/r$
constant when sweeping $r$ maintains consistent update magnitude.  After training,
$\mB\mA$ merges into $\mW_0$ for zero inference overhead.

\redflags
\begin{itemize}
    \item Cannot state the formula $\mW = \mW_0 + \mB\mA$
    \item Does not know that $\mB$ is initialized to zero (and why)
    \item Cannot explain the intrinsic dimensionality argument for why low-rank works
    \item Unaware of the merging property for zero inference overhead
\end{itemize}

\interviewtest{Whether the candidate understands LoRA at a mathematical level, not just
as a library call.  Can they explain the low-rank hypothesis, the initialization choice,
and the practical implications (merging, multi-task serving)?}

\goodgreat
\begin{itemize}
    \item \badgeLfive{L5} States the formula, explains the initialization and merging
          property, gives reasonable defaults for $r$ and $\alpha$
    \item \badgeLsix{L6} Derives the parameter count savings quantitatively, explains
          the intrinsic dimensionality argument, discusses which modules to apply LoRA
          to and why, and compares with other PEFT methods
    \item \badgeLseven{L7} Discusses the spectral analysis of weight updates, connects
          to matrix completion theory, addresses when LoRA fails (high-rank updates
          needed for large distribution shifts), and describes multi-adapter serving
          architectures
\end{itemize}

\followups{How does the gradient flow through frozen weights in LoRA?  What is
QLoRA?  When would LoRA fail?  How would you serve 100 different LoRA adapters
from one base model?}
\end{interviewq}

\bigskip

%% QUESTION 2 %%
\begin{interviewq}
\textbf{Q: When would you use LoRA vs.\ full fine-tuning? When would LoRA fail?}

\badgeLfive{L5} \quad \badgetype{Trade-off} \quad \badgefreq{\freqhigh}

\stronganswer
\textbf{Use LoRA when:} (1) The base model is large (7B+) and full fine-tuning
is memory-prohibitive. (2) You need to serve multiple tasks from one base model.
(3) The downstream task is reasonably close to the pre-training distribution.
(4) You need fast iteration (LoRA trains 2--3$\times$ faster due to fewer
optimizer states).

\textbf{Use full fine-tuning when:} (1) Maximum quality is required and you have
the compute budget. (2) The domain shift is extreme (e.g., pre-trained on English
web text, fine-tuning for specialized medical code in another language). (3) The
model is small enough that full fine-tuning is feasible (under 1B parameters).
(4) You are doing continued pre-training (not just task adaptation).

\textbf{LoRA fails when:} (1) The required weight update is genuinely high-rank---
this happens when the target distribution differs substantially from pre-training
(the model needs to ``unlearn'' and ``relearn'' many patterns). (2) The task requires
modifying the embedding layer significantly (e.g., adding a new language). (3) Very
low rank is used for a complex task---the bottleneck starves the model of capacity.

\redflags
\begin{itemize}
    \item Says ``always use LoRA'' or ``always use full fine-tuning''
    \item Cannot articulate a concrete scenario where LoRA would fail
    \item Does not mention memory/compute as the primary practical motivation
\end{itemize}

\interviewtest{Whether the candidate can make nuanced engineering decisions about
adaptation methods rather than defaulting to one approach.  This tests production
maturity and understanding of tradeoffs.}

\goodgreat
\begin{itemize}
    \item \badgeLfive{L5} Gives clear criteria for each approach with concrete examples
    \item \badgeLsix{L6} Quantifies the memory and compute savings, discusses the rank
          vs.\ quality tradeoff with specific numbers, mentions QLoRA as a middle ground
    \item \badgeLseven{L7} Discusses the theoretical conditions under which low-rank
          approximation is insufficient, connects to the singular value distribution of
          gradient updates, and describes hybrid approaches (full fine-tuning some layers,
          LoRA for others)
\end{itemize}

\followups{How would you decide the rank for a specific task?  What if you need to
add a new language to the model?  How does LoRA interact with quantization?}
\end{interviewq}

\bigskip

%% QUESTION 3 %%
\begin{interviewq}
\textbf{Q: Compare LoRA, prefix tuning, and prompt tuning. What are the tradeoffs?}

\badgeLsix{L6} \quad \badgetype{Conceptual} \quad \badgefreq{\freqmed}

\stronganswer
\textbf{LoRA} adds low-rank matrices to attention (and optionally FFN) weights.
Trainable params: 0.1--1\%.  After training, merges into base weights for zero
inference overhead.  Quality: 90--99\% of full fine-tuning.  Best all-around
PEFT method.

\textbf{Prefix tuning} prepends learnable virtual tokens to keys and values at
every layer.  Trainable params: $\sim$0.1\%.  Adds inference overhead (extra KV
per layer) and consumes context window positions.  Quality: 85--95\%.  Cannot
merge into base weights.

\textbf{Prompt tuning} adds learnable soft tokens to the input embedding only (not
every layer).  Trainable params: $< 0.1\%$.  Simplest method, fewest parameters.
Quality scales with model size: poor below 1B, competitive above 10B.

\textbf{Key tradeoffs:} (1) LoRA is the only method with zero inference overhead
(via merging). (2) Prompt tuning is the simplest but works well only for large
models. (3) Prefix tuning is more expressive than prompt tuning (touches every
layer) but less so than LoRA (cannot modify weight behavior, only prepend context).
(4) LoRA modifies \emph{how} the model processes tokens; prefix/prompt tuning
modifies \emph{what} the model processes (by adding virtual context).

\redflags
\begin{itemize}
    \item Conflates prefix tuning with prompt tuning
    \item Does not mention LoRA's merging advantage
    \item Cannot explain the architectural difference (weight modification vs.\ input modification)
\end{itemize}

\interviewtest{Whether the candidate understands the \emph{mechanism} of each method,
not just their names.  The weight vs.\ input modification distinction is the key
conceptual insight.}

\goodgreat
\begin{itemize}
    \item \badgeLfive{L5} Correctly describes all three methods with parameter counts
          and quality ranges
    \item \badgeLsix{L6} Explains the architectural distinction (weight modification
          vs.\ virtual token injection), discusses when each is appropriate, and mentions
          adapters and IA3 as additional alternatives
    \item \badgeLseven{L7} Discusses the theoretical expressiveness hierarchy, connects
          prefix tuning to the kernel perspective on attention, and analyzes how
          different methods compose (e.g., LoRA + prompt tuning)
\end{itemize}

\followups{Can you combine LoRA with prefix tuning?  How does prompt tuning relate
to discrete prompt engineering?  What is IA3?}
\end{interviewq}

\bigskip

%% QUESTION 4 %%
\begin{interviewq}
\textbf{Q: How would you fine-tune an LLM for a domain-specific task with only
1,000 labeled examples?}

\badgeLfive{L5} \quad \badgetype{System Design} \quad \badgefreq{\freqhigh}

\stronganswer
With only 1,000 examples, full fine-tuning of a large LLM risks severe overfitting.
My approach:

(1)~\textbf{Start with LoRA} (rank 16--32) on a strong base model (e.g., LLaMA-2 7B
or 13B).  Low-rank adaptation with few parameters is naturally regularizing.
(2)~\textbf{Data quality audit:} ensure the 1,000 examples are diverse, correct,
and representative.  1,000 high-quality examples can be remarkably effective (LIMA
demonstrated this).
(3)~\textbf{Low learning rate} ($1\text{e-}5$ to $2\text{e-}5$) with cosine
schedule and warmup.
(4)~\textbf{Train for 3--5 epochs} with early stopping based on validation loss.
(5)~\textbf{Data augmentation:} use the base model itself to generate synthetic
examples (bootstrap from the 1,000), filter by quality, and add to training.
(6)~\textbf{Evaluation:} hold out 10--15\% for validation; use task-specific
metrics, not just loss.

If quality is still insufficient: (7)~continued pre-training on domain-specific
unlabeled text (before SFT) to adapt the model's representations to the domain.
(8)~Few-shot prompting as a baseline---for some tasks, a well-prompted LLM with
5 examples outperforms fine-tuning on 1,000 examples when the fine-tuning data
is noisy.

\redflags
\begin{itemize}
    \item Proposes full fine-tuning of a 70B model on 1,000 examples without
          acknowledging overfitting risk
    \item Does not mention LoRA or any PEFT method
    \item Ignores data quality and diversity considerations
    \item Does not mention evaluation strategy
\end{itemize}

\interviewtest{Practical ML engineering judgment.  Can the candidate design an
end-to-end fine-tuning strategy with appropriate regularization, data management,
and evaluation under real-world constraints?}

\goodgreat
\begin{itemize}
    \item \badgeLfive{L5} Proposes LoRA with reasonable hyperparameters, mentions
          overfitting risk and early stopping
    \item \badgeLsix{L6} Discusses continued pre-training as a preliminary step,
          synthetic data augmentation, compares with few-shot prompting baseline,
          and designs an evaluation strategy
    \item \badgeLseven{L7} Designs a full pipeline: domain-adaptive pre-training,
          data augmentation with self-consistency filtering, LoRA fine-tuning with
          hyperparameter search over rank, active learning loop to acquire more
          labeled data, and production monitoring for domain drift
\end{itemize}

\followups{How would you do continued pre-training?  How would you decide
between LoRA and few-shot prompting?  What if the 1,000 examples are noisy?}
\end{interviewq}

\bigskip

%% QUESTION 5 %%
\begin{interviewq}
\textbf{Q: What is catastrophic forgetting? How do you mitigate it during fine-tuning?}

\badgeLfive{L5} \quad \badgetype{Conceptual} \quad \badgefreq{\freqhigh}

\stronganswer
Catastrophic forgetting occurs when fine-tuning on a downstream task causes the model
to lose knowledge acquired during pre-training.  The parameters shift to optimize
the new loss, destructively overwriting representations that encoded general
language understanding.  The smaller the fine-tuning dataset and the higher the
learning rate, the more severe the forgetting.

Mitigations: (1)~\textbf{Low learning rate} ($1\text{e-}5$ to $5\text{e-}5$)---the
single most important factor. (2)~\textbf{Gradual unfreezing} (ULMFiT)---train the
task head first, then unfreeze layers top-to-bottom. (3)~\textbf{Regularization}---
weight decay keeps parameters near pre-trained values; EWC penalizes changes to
parameters important for pre-training. (4)~\textbf{Early stopping}---fine-tuning
peaks in 2--3 epochs. (5)~\textbf{Data replay}---mix pre-training-style data into
fine-tuning batches. (6)~\textbf{PEFT methods} (LoRA, adapters)---by only modifying
a small fraction of parameters, the vast majority of pre-trained knowledge is
preserved by design.

PEFT methods are arguably the most effective mitigation because they structurally
prevent forgetting: the base weights are literally frozen.

\redflags
\begin{itemize}
    \item Cannot define catastrophic forgetting
    \item Only mentions ``use a small learning rate'' with no other mitigations
    \item Does not connect PEFT methods to forgetting prevention
\end{itemize}

\interviewtest{Understanding of a fundamental challenge in transfer learning and the
ability to propose multiple, principled mitigations.}

\goodgreat
\begin{itemize}
    \item \badgeLfive{L5} Defines forgetting clearly, lists 3+ mitigations with
          rationale for each
    \item \badgeLsix{L6} Discusses EWC formally, connects to continual learning
          literature, explains why PEFT methods structurally prevent forgetting
    \item \badgeLseven{L7} Discusses the loss landscape geometry of forgetting
          (mode connectivity, linear mode connectivity), analyzes the tradeoff
          between plasticity and stability, and connects to multi-task learning
\end{itemize}

\followups{What is Elastic Weight Consolidation?  How does forgetting differ for
encoder-only vs.\ decoder-only models?  Can you recover from forgetting?}
\end{interviewq}

\bigskip

%% QUESTION 6 %%
\begin{interviewq}
\textbf{Q: Why is data quality more important than data quantity for SFT?}

\badgeLfive{L5} \quad \badgetype{Conceptual} \quad \badgefreq{\freqmed}

\stronganswer
During pre-training, the model already acquires broad language understanding and
factual knowledge from trillions of tokens.  SFT does not teach the model
\emph{new knowledge}---it teaches the model a \emph{behavioral format}: how to
structure responses, follow instructions, and be helpful.  The LIMA paper (Zhou
et al., 2023) showed that 1,000 high-quality examples are sufficient to align a
65B model's behavior, while Alpaca-style models trained on 52K noisy synthetic
examples performed worse on many dimensions.

The reason: noisy or inconsistent SFT data sends conflicting signals about the
desired behavior, causing the model to learn an averaged (mediocre) response
style.  High-quality data provides a clear, coherent target.  Additionally,
because SFT datasets are small relative to pre-training data, each example has
outsized influence on the model's behavior---one bad example in 1,000 has far
more impact than one bad document in a trillion-token corpus.

This principle also extends to diversity: covering many task types matters more
than having many examples of the same type.  10 diverse, high-quality examples
across 100 task types teach more than 1,000 examples of a single task type.

\redflags
\begin{itemize}
    \item Says ``more data is always better'' without qualification
    \item Does not distinguish the role of SFT from pre-training
    \item Unaware of the LIMA result or similar findings
\end{itemize}

\interviewtest{Whether the candidate understands the distinct roles of pre-training
(knowledge acquisition) and SFT (behavior formatting), and can reason about data
efficiency.}

\goodgreat
\begin{itemize}
    \item \badgeLfive{L5} Explains the knowledge-vs.-behavior distinction and
          references LIMA
    \item \badgeLsix{L6} Discusses the data quality dimensions (correctness, diversity,
          consistency, complexity), contrasts with pre-training data strategy, and
          explains why synthetic data quality is critical for SFT
    \item \badgeLseven{L7} Connects to the broader alignment tax concept, discusses
          how data quality interacts with model scale, and analyzes the failure modes
          of low-quality SFT data (sycophancy, hallucination amplification,
          inconsistent safety behavior)
\end{itemize}

\followups{How would you create high-quality SFT data?  What role does synthetic
data play?  How do you evaluate SFT data quality before training?}
\end{interviewq}

\bigskip

%% QUESTION 7 %%
\begin{interviewq}
\textbf{Q: Explain ELECTRA's pre-training approach. Why is it more efficient than MLM?}

\badgeLsix{L6} \quad \badgetype{Conceptual} \quad \badgefreq{\freqmed}

\stronganswer
ELECTRA replaces masked language modeling with replaced token detection (RTD).
A small generator network (trained with MLM) produces plausible replacement tokens
for masked positions.  The main model (the discriminator) receives the corrupted
sequence and classifies each token as ``real'' (original) or ``replaced'' (from the
generator).  The generator is discarded after pre-training; only the discriminator
is used for fine-tuning.

The efficiency advantage is dramatic: MLM provides training signal at only the 15\%
of tokens that are masked.  ELECTRA provides signal at \emph{100\%} of tokens---
every position is classified.  This $\sim$7$\times$ increase in per-example signal
means ELECTRA matches BERT's performance with roughly 1/4 the compute.  With equal
compute, ELECTRA significantly outperforms BERT.

The key insight is that detection is an easier task than generation (binary
classification vs.\ predicting from a 30K+ vocabulary), so even a small discriminator
can learn effectively from every position.  The generator only needs to be good
enough to produce ``hard'' negatives---tokens that are plausible but wrong.

\redflags
\begin{itemize}
    \item Confuses the generator and discriminator roles
    \item Cannot explain why 100\% vs.\ 15\% signal matters for efficiency
    \item Describes ELECTRA as a GAN (it is not---the generator is trained with MLM,
          not adversarially)
\end{itemize}

\interviewtest{Whether the candidate understands alternative pre-training objectives
beyond the standard BERT/GPT approaches, and can reason about sample efficiency.}

\goodgreat
\begin{itemize}
    \item \badgeLfive{L5} Correctly describes the generator-discriminator setup and
          the efficiency argument
    \item \badgeLsix{L6} Explains why ELECTRA is not a GAN (different training
          objectives), discusses the generator size tradeoff (too small: easy negatives;
          too large: wastes compute), and compares quantitatively with MLM
    \item \badgeLseven{L7} Discusses the connection to noise-contrastive estimation,
          analyzes the theoretical sample complexity advantage, and connects to
          energy-based models and discriminative vs.\ generative pre-training
\end{itemize}

\followups{Why is ELECTRA not a GAN?  What happens if the generator is too good?
How does ELECTRA compare to T5's span corruption?}
\end{interviewq}

\bigskip

%% QUESTION 8 %%
\begin{interviewq}
\textbf{Q: What is the difference between feature extraction and fine-tuning?
When would you use each?}

\badgeLfive{L5} \quad \badgetype{Conceptual} \quad \badgefreq{\freqmed}

\stronganswer
\textbf{Feature extraction} freezes all pre-trained weights and trains only a
task-specific head (typically a linear layer or small MLP) on top of the model's
output representations.  The pre-trained model is a static feature extractor.
\textbf{Fine-tuning} updates all (or most) model parameters end-to-end on the
downstream task, allowing the representations themselves to adapt.

Fine-tuning typically outperforms feature extraction by 2--5\% absolute because
it adapts the internal representations to the specific task distribution.  However,
feature extraction is preferred when: (1)~labeled data is extremely limited
(tens of examples---fine-tuning would overfit). (2)~Compute is constrained
(feature extraction requires only one forward pass to extract features, then
trains a lightweight classifier). (3)~You need to preserve the base model
exactly (e.g., using the same model for multiple tasks without any modification).
(4)~As a quick baseline before investing in fine-tuning.

In practice, PEFT methods like LoRA occupy the middle ground: they update more
than feature extraction (the low-rank matrices adapt the representations) but far
less than full fine-tuning, combining the quality advantage of adaptation with the
efficiency of limited parameter updates.

\redflags
\begin{itemize}
    \item Cannot articulate the difference clearly
    \item Says fine-tuning is always better (ignores low-data and compute-constrained
          scenarios)
    \item Does not mention PEFT as a middle ground
\end{itemize}

\interviewtest{Whether the candidate understands the spectrum of adaptation strategies
and can make appropriate choices based on constraints.}

\goodgreat
\begin{itemize}
    \item \badgeLfive{L5} Correctly distinguishes the two approaches, gives conditions
          for choosing each
    \item \badgeLsix{L6} Discusses layer-wise feature quality (earlier layers more
          general, later layers more task-specific), mentions probing experiments, and
          positions PEFT methods on the spectrum
    \item \badgeLseven{L7} Discusses the representation learning perspective (when do
          pre-trained features need adaptation vs.\ when are they sufficient), connects
          to the lottery ticket hypothesis, and analyzes how model scale affects the
          feature-extraction-vs.-fine-tuning tradeoff
\end{itemize}

\followups{Which layer's representations are best for feature extraction?  How does
model size affect the gap between feature extraction and fine-tuning?  What is
probing and why does it matter?}
\end{interviewq}

\bigskip

%% QUESTION 9 %%
\begin{interviewq}
\textbf{Q: Walk through the QLoRA paper. How does it enable fine-tuning a 65B
model on a single GPU?}

\badgeLsix{L6} \quad \badgetype{Mathematical} \quad \badgefreq{\freqmed}

\stronganswer
QLoRA combines three innovations to make large-model fine-tuning memory-feasible.

(1)~\textbf{NF4 quantization:} Base model weights are stored in 4-bit precision
using a data type whose quantization bins are optimally spaced for normally
distributed values (which pre-trained weights empirically follow).  This reduces
the base model from $\sim$130\,GB (FP16) to $\sim$33\,GB for a 65B model.

(2)~\textbf{Double quantization:} Each block of 64 weights has a FP32 quantization
constant (scale factor).  These constants are themselves quantized to FP8, reducing
the overhead from $\sim$0.5 bits/parameter to $\sim$0.13 bits/parameter---saving
$\sim$3\,GB for a 65B model.

(3)~\textbf{Paged optimizers:} When gradient computation causes GPU memory spikes,
optimizer states are offloaded to CPU RAM via CUDA unified memory, then paged back
when needed.

During the forward pass, 4-bit weights are dequantized to BF16 for computation.
LoRA adapters ($\mA, \mB$) remain in BF16 throughout.  Gradients only update the
LoRA parameters---the quantized base weights are never modified.  Total memory
for a 65B model: $\sim$33\,GB base + $\sim$0.5\,GB LoRA + optimizer states for
LoRA only $\approx$ fits in a single 48\,GB A6000 GPU.

\redflags
\begin{itemize}
    \item Does not know what NF4 is or why normal-float is better than uniform INT4
    \item Cannot explain double quantization
    \item Thinks the 4-bit weights are trained (they are frozen)
\end{itemize}

\interviewtest{Whether the candidate understands the systems-level innovations that
make large-scale PEFT practical, not just the algorithmic idea.}

\goodgreat
\begin{itemize}
    \item \badgeLfive{L5} Describes the three innovations and the overall memory savings
    \item \badgeLsix{L6} Explains NF4 vs.\ uniform quantization, quantifies memory at
          each stage, and discusses the quality-memory tradeoff
    \item \badgeLseven{L7} Discusses the interaction between quantization error and LoRA
          rank, compares with GPTQ/AWQ post-training quantization, and analyzes when
          QLoRA quality degrades (very low bit-widths, models with outlier activations)
\end{itemize}

\followups{What is the difference between QLoRA and GPTQ?  How does quantization
affect LoRA's effective rank?  Can you do QLoRA with 2-bit quantization?}
\end{interviewq}

\bigskip

%% QUESTION 10 %%
\begin{interviewq}
\textbf{Q: Compare the major pre-training objectives: MLM, CLM, span corruption,
and ELECTRA. When would you choose each?}

\badgeLsix{L6} \quad \badgetype{Trade-off} \quad \badgefreq{\freqmed}

\stronganswer
\textbf{MLM (BERT):} Bidirectional, predicts 15\% of tokens.  Best for
understanding tasks (classification, NER, retrieval) where full context is
available at inference.  Limitation: 15\% signal is sample-inefficient; cannot
generate text natively.

\textbf{CLM (GPT):} Unidirectional, predicts every token left-to-right.  Best for
generation, in-context learning, and general-purpose LLMs.  100\% signal makes it
more data-efficient.  Limitation: no right context for understanding tasks.

\textbf{Span corruption (T5):} Bidirectional encoder + autoregressive decoder.  Best
for structured input-output tasks (translation, summarization, structured extraction).
The encoder-decoder split naturally separates understanding and generation.

\textbf{ELECTRA (RTD):} Bidirectional, 100\% signal via replaced token detection.
Most sample-efficient for encoder models.  Best when compute is limited and you
need a strong encoder for understanding tasks.

Decision framework: (1)~If generation is the primary use case $\rightarrow$ CLM.
(2)~If understanding/classification with unlimited compute $\rightarrow$ ELECTRA.
(3)~If structured seq2seq tasks $\rightarrow$ span corruption. (4)~If a simple,
well-understood baseline $\rightarrow$ MLM.  In practice, CLM dominates the LLM
era because decoder-only models handle both generation and understanding (via
in-context learning), even though encoder-only models may be superior for
pure understanding tasks.

\redflags
\begin{itemize}
    \item Can only describe one or two objectives
    \item Does not mention the signal-per-token efficiency difference
    \item Cannot articulate when to choose each objective
\end{itemize}

\interviewtest{Breadth and depth of understanding across the pre-training landscape,
and the ability to reason about tradeoffs rather than memorize facts.}

\goodgreat
\begin{itemize}
    \item \badgeLfive{L5} Correctly describes all four objectives and gives use cases
    \item \badgeLsix{L6} Compares sample efficiency quantitatively, discusses UL2
          (combining multiple objectives) and PrefixLM, and explains why CLM dominates
          despite MLM's theoretical advantage for understanding
    \item \badgeLseven{L7} Discusses the information-theoretic perspective (mutual
          information between masked/predicted tokens and context), analyzes the
          pretraining-objective-to-capability mapping, and connects to the broader
          question of whether architectural inductive biases or objective functions
          matter more at scale
\end{itemize}

\followups{What is UL2?  Why did the field converge on CLM despite MLM being better
for understanding?  How does denoising (BART) compare to span corruption (T5)?}
\end{interviewq}


% ============================================================================
% CHAPTER SUMMARY
% ============================================================================
\section*{Chapter Summary}

\begin{enumerate}
    \item \textbf{Pre-training objectives} determine model capabilities:
          MLM provides bidirectional understanding, CLM enables generation and
          in-context learning, span corruption bridges both via encoder-decoder
          architecture, and ELECTRA achieves maximum sample efficiency through
          replaced token detection.

    \item \textbf{Pre-training data quality} is the most underappreciated
          factor in LLM development.  The pipeline of language ID, deduplication,
          quality filtering, and domain mixing directly determines what the model
          knows and how well it performs.

    \item \textbf{Transfer learning} works because pre-training learns
          universal language representations---syntax, semantics, world
          knowledge---that transfer across downstream tasks with minimal
          labeled data.

    \item \textbf{Full fine-tuning} updates all parameters and achieves the
          highest quality, but risks catastrophic forgetting.  Mitigations
          include low learning rates, gradual unfreezing, regularization, and
          early stopping.

    \item \textbf{LoRA} is the dominant PEFT method: it adds low-rank matrices
          $\mB\mA$ to frozen weights, trains 0.1--1\% of parameters, and merges
          for zero inference overhead.  The intrinsic dimensionality of
          fine-tuning justifies the low-rank assumption.

    \item \textbf{QLoRA} extends LoRA to models that would not otherwise fit
          in memory by quantizing the base model to 4-bit (NF4) with double
          quantization, enabling 65B-parameter fine-tuning on a single GPU.

    \item \textbf{Other PEFT methods}---prefix tuning, prompt tuning, adapters,
          IA3, DoRA---offer different tradeoffs in parameter efficiency,
          inference overhead, and quality.  LoRA dominates in practice due to
          its combination of quality, efficiency, and zero inference overhead.

    \item \textbf{Supervised fine-tuning (SFT)} transforms pre-trained models
          into instruction-following assistants.  Data quality matters far more
          than quantity: 1,000 high-quality examples can be sufficient (LIMA).
          Loss masking on assistant tokens only is a critical implementation
          detail.

    \item \textbf{The three-stage pipeline}---pre-training, SFT, alignment---is
          the standard recipe for building production LLMs.  Each stage serves a
          distinct purpose: knowledge acquisition, behavioral formatting, and
          quality refinement.
\end{enumerate}
